\begin{table*}[t]
\centering
\scriptsize
\renewcommand{\arraystretch}{0.9}
\begin{tabularx}{\linewidth}{>{\raggedright\arraybackslash}p{0.18\linewidth}>{\raggedright\arraybackslash}p{0.22\linewidth}>{\raggedright\arraybackslash}p{0.25\linewidth}X}
\toprule
\textbf{Lineage} & \textbf{Classical question} & \textbf{RAG-specific version} & \textbf{Why it matters for robust RAG} \\
\midrule
Selective prediction & When should a model abstain to control risk? & Abstain when retrieved evidence is insufficient, conflicting, stale, or corrupted. & Selective risk must condition on evidence state, not only answer confidence. \\
Calibration & Does confidence match correctness? & Does confidence match answer correctness, evidence sufficiency, citation support, and retrieval usefulness? & A calibrated answer probability can still hide unreliable evidence. \\
Active clarification & When should the system ask a user for more information? & Ask clarification before retrieval when the query is ambiguous or underspecified. & Clarification can be the correct action even when retrieval could produce a plausible answer. \\
Cost-sensitive retrieval & When is additional computation worth it? & Decide whether to retrieve, re-retrieve, verify, rerank, or escalate under latency/cost budgets. & Robustness gains should be reported with action cost and coverage, not only accuracy. \\
Decision-theoretic control & Which action minimizes expected risk? & Choose among answer, abstain, expose conflict, cite uncertainty, retrieve again, or quarantine evidence. & Action quality is the bridge between evidence reliability estimates and user-facing behavior. \\
\bottomrule
\end{tabularx}
\caption{Action-control lineage. The framework imports ideas from selective prediction, calibration, clarification, and cost-sensitive decision-making, but makes retrieved evidence state the object being controlled.}
\label{tab:action-control-lineage}
\end{table*}

